\begin{abstract}
Although large-scale unsupervised language models (LMs) acquire extensive world knowledge and certain reasoning abilities, exerting fine-grained control over their outputs remains challenging due to the inherently unsupervised process by which they are trained. Current approaches to increase model steerability generally rely on obtaining human-provided preference labels comparing the quality of model outputs and subsequently adapting the LM to better reflect these preferences, frequently via reinforcement learning from human feedback (RLHF). Nevertheless, RLHF is often a complicated and unstable pipeline that involves first constructing a reward model representing human judgment, then employing reinforcement learning to adjust the unsupervised LM towards optimizing this learned reward function, all while maintaining closeness to the original model parameters. In this study, we propose a novel parameterization for the RLHF reward model that admits a closed-form solution for the associated optimal policy, which, in turn, allows the canonical RLHF task to be addressed using nothing more than a straightforward classification objective. We term this technique \textit{Direct Preference Optimization} (DPO). DPO offers stable performance, strong empirical results, and minimal computational overhead, obviating the necessity for generating samples from the LM during fine-tuning or for extensive hyperparameter tuning. Our experimental results demonstrate that DPO is able to fine-tune LMs in alignment with human preferences as well as—or even surpassing—current leading approaches. Significantly, DPO-based fine-tuning not only outperforms PPO-driven RLHF in steering generation sentiment but also achieves comparable or superior outcomes in summarization and single-turn dialogue, all while greatly reducing the complexity of implementation and training. 
\end{abstract}

\section{Introduction}
Large-scale unsupervised language models (LMs), trained on immense corpora, have demonstrated unexpected emergent abilities~\citep{chowdhery2022palm, brown2020language, touvron2023llama,bubeck2023sparks}. Nonetheless, their training data, authored by humans with diverse intentions, competencies, and backgrounds, includes a wide range of behaviors—some of which may be undesirable to replicate. For instance, it is beneficial if an AI-based coding assistant can \textit{recognize} typical programming errors to facilitate corrections, yet, during code generation, the system should be oriented toward exemplifying the highest levels of coding proficiency found in the training material, even if such examples are infrequent. In another case, it may be useful for a language model to \textit{know} about a misconception commonly accepted by half the population, but it is undesirable for the model to state the misconception as fact in 50\% of related outputs. Thus, it is vital to distinguish between a model's extensive range of \textit{knowledge and skills} and the specific \emph{behaviors and responses} we desire it to exhibit, to ensure that AI models are reliable, effective, and manageable \citep{ouyang2022training}. Existing strategies often employ reinforcement learning (RL) to adjust LMs according to human preference data. In this work, we demonstrate that the RL-based objective in these approaches can actually be realized precisely through a straightforward binary cross-entropy loss, which considerably streamlines the process of learning from preferences.

Broadly, current approaches impart preferred behavior into language models using carefully constructed datasets of human judgments, embodying qualities regarded as safe and useful by users. This phase of preference modeling is typically carried out following an extensive phase of unsupervised pre-training. While supervised fine-tuning on exemplary human-generated responses is a direct method, reinforcement learning from human or AI feedback (RLHF/RLAIF; \citep{christiano2017deep,bai2022constitutional}) has achieved the most impressive results. In RLHF, a reward function is learned based on human preference data, and the language model's output distribution is then shaped via RL to favor responses with higher reward, all while constraining divergence from the pretrained model. Although RLHF techniques yield models with remarkable abilities in dialogue and programming, their training pipelines are considerably more complicated compared to supervised approaches, requiring multiple language models, as well as repeated policy sampling during training, and thus lead to substantial computational demands.

This work demonstrates a method for optimizing language models in accordance with human preferences, bypassing the need for explicit reward modeling or reinforcement learning techniques. We introduce \textit{Direct Preference Optimization} (DPO), a procedure that implicitly targets the same reward-maximizing objective under a KL-divergence regularization as traditional RLHF algorithms, while offering a more straightforward implementation and training process. DPO intuitively adjusts the model by boosting the log-likelihood ratio between preferred and non-preferred outputs, yet employs an adaptive, sample-specific weighting mechanism to avoid the collapse that may occur when employing a naive ratio-based objective. As with earlier algorithms, DPO leverages a theoretical model of preferences (e.g., the Bradley-Terry model; \cite{bradley1952rankanalysis}) to quantify the alignment between a proposed reward function and observed preference data. Distinctly, instead of first fitting a reward model using this preference model and subsequently training a policy to maximize it, DPO utilizes a variable substitution that expresses the preference loss directly in terms of the policy itself. This allows DPO to directly update the policy using a basic binary cross entropy loss derived from datasets containing human rankings of responses, resulting in an optimal policy with respect to an implicitly defined reward shaped by the observed preferences.

The principal contribution of this work is the presentation of Direct Preference Optimization (DPO), an uncomplicated algorithm for preference-based language model training that does not rely on reinforcement learning. Through empirical studies, we find that DPO achieves comparable, if not superior, performance to prevailing approaches such as PPO-based RLHF on preference-driven benchmarks including sentiment control, text summarization, and dialogue generation, scaling to models containing up to 6 billion parameters.

\section{Related Work}

As self-supervised language models scale up, they increasingly demonstrate the ability to perform certain tasks with no additional training (zero-shot) \citep{radford2019language} or with minimal prompting (few-shot) \citep{gpt3,megatron,chowdhery2022palm}. Nonetheless, augmenting their capabilities and better aligning them with user intent often requires fine-tuning on curated collections of instructions paired with human-generated completions \citep{mishra-etal-2022-cross,sanh2022multitask,chung2022scaling,thoppilan2022lamda}. This process, known as `instruction-tuning,' not only boosts the models' ability to execute a wider range of instructions beyond those seen during fine-tuning but also enhances their general practicality \citep{chung2022scaling}. Although instruction tuning has proved effective, collecting \textit{relative} human evaluations comparing response quality is frequently less demanding than obtaining expert-level demonstrations. As a result, recent research has shifted towards fine-tuning LLMs using preference data derived from such comparative human assessments, which has led to measurable gains in areas including translation \citep{kreutzer-etal-2018-reliability}, summarization \citep{stiennon2022learning,ziegler2020finetuning}, narrative generation \citep{ziegler2020finetuning}, and following instructions \citep{ouyang2022training,ramamurthy2023is}. The standard methodology involves initially training a reward model—guided by a preference framework like the Bradley-Terry model \citep{bradley1952rankanalysis}—to fit human preference data, after which the language model is optimized via reinforcement learning approaches such as REINFORCE \citep{williams1992reinforce}, proximal policy optimization (PPO; \cite{schulman2017proximal}), or related algorithms \citep{ramamurthy2023is}. Another closely related direction explores the generation of additional synthetic preference data—targeted at features like safety or non-harmfulness—using LLMs that have already been instruction-tuned with human feedback; this strategy relies mainly on limited human oversight provided through textual rubrics for model-generated annotations \citep{bai2022constitutional}. Collectively, these advances reflect an intersection between prior efforts to train language models via reinforcement learning to meet diverse objectives~\citep{Ranzato2015SequenceLT,paulus2018a,wu2018learning} and broader strategies for leveraging human preference information \citep{christiano2017deep,kupcsik2018learning}. Despite the efficiency and accessibility of relative human preference data, reinforcement learning-based fine-tuning of large language models introduces substantial practical obstacles; the current work presents a theoretically grounded alternative for optimizing with relative preferences that does not rely on reinforcement learning.

Beyond linguistic contexts, the study of policy learning from preferences has been explored within both bandit and reinforcement learning paradigms, giving rise to a variety of methodological innovations. When the feedback consists of action preferences or rankings instead of explicit rewards, the problem is framed as a contextual dueling bandit (CDB; \cite{yue2012karmed,dudik2015contextual}). In CDBs, lacking absolute reward signals, theoretical guarantees rely on the notion of a \textit{von Neumann winner}—a policy that is expected to outperform any alternative policy at least half the time \citep{dudik2015contextual}. Notably, CDB approaches collect preference labels in an online manner, in contrast with the offline, batch-wise setting typically assumed when learning from human preference judgments, where fixed datasets of annotated action pairs are employed \citep{yan2022human}. Likewise, in \textit{preference-based RL} (PbRL), the agent relies on binary preferences produced by an unspecified scoring mechanism in lieu of standard reward signals \citep{BusaFekete2014,ruiz2023dueling}. PbRL solutions span a range of techniques, often leveraging off-policy data, but generally proceed by first inferring the latent reward or scoring model and then optimizing with respect to this inferred function \citep{jain2013learning,BusaFekete2014,christiano2017deep,sadigh2017active,kupcsik2018learning}. By contrast, we advocate for a unified policy learning strategy that eschews the intermediate reward model estimation and directly optimizes policies to adhere to observed preferences.

\section{Preliminaries}\label{section:prelims}

We provide an overview of the RLHF process as formulated by \citeauthor{ziegler2020finetuning}, with subsequent extensions by \citep{stiennon2022learning, bai2022training, ouyang2022training}. The workflow typically consists of three stages: 1) supervised fine-tuning (SFT); 2) reward learning from preference data; and 3) reinforcement learning optimization.

\textbf{SFT}: The RLHF procedure is generally initiated by adapting a pre-trained language model through supervised learning, using carefully curated datasets for specific downstream objectives (such as dialogue or summarization). This produces a model denoted $\pi^\text{SFT}$.

\textbf{Reward Modelling Phase}: In this stage, prompts $x$ are input to the SFT model to generate pairs of candidate completions $(y_1, y_2)\sim \pi^\text{SFT}(y \mid x)$. Human annotators then evaluate these pairs and express a preference, represented as $y_w\succ y_l \mid x$, where $y_w$ and $y_l$ are, respectively, the favored and less favored completions from the pair. The preferences are postulated to arise from a hidden reward function $r^*(y, x)$ that is not directly observed. Several preference modeling techniques exist, with the Bradley-Terry (BT) \cite{bradley1952rankanalysis} model being widely adopted (though more general approaches such as the Plackett-Luce model \citep{plackett1975analysis, luce2012individual} can also be utilized if multiple rankings are available). Under the BT framework, the true preference distribution $p^*$ is formulated as follows:

\begin{equation}\label{eq:bradley-terry}
    p^*(y_1\succ y_2 \mid x)=\frac{\exp\left(r^*(x, y_1)\right)}{\exp\left(r^*(x, y_1)\right) + \exp\left(r^*(x, y_2)\right)}.
\end{equation}

Suppose we have a fixed dataset of preference comparisons $\mathcal{D}=\bigl\{x^{(i)}, y_w^{(i)}, y_l^{(i)}\bigr\}_{i=1}^N$, drawn according to $p^*$. One can then define a reward model $r_{\phi}(x, y)$ with parameters $\phi$ and fit these parameters by maximizing the likelihood. Recasting this as a binary classification problem yields the following negative log-likelihood objective:

\begin{equation}\label{eq:reward_model}
    \mathcal{L}_R(r_{\phi}, \mathcal{D}) = -\mathbb{E}_{(x, y_w, y_l)\sim \mathcal{D}}\bigl[\log \sigma(r_{\phi}(x, y_w)- r_{\phi}(x, y_l))\bigr]
\end{equation}

with $\sigma$ denoting the logistic (sigmoid) function. When training LMs, $r_{\phi}(x, y)$ is typically initialized using the SFT model $\pi^\text{SFT}(y \mid x)$, supplemented by a linear head placed atop the final transformer output to generate a scalar reward estimate \cite{ziegler2020finetuning}. In order to promote stability and reduce reward variance, prior work applies normalization, enforcing $\mathbb{E}_{x,y\sim \mathcal{D}}\left[r_\phi(x, y)\right] = 0$ for each $x$.

\textbf{RL Fine-Tuning Phase}: In the subsequent reinforcement learning stage, the reward model provides learning signals to the language model. Consistent with prior approaches~\citep{jaques2017sequence, jaques2020human}, the objective is expressed as

\begin{equation}\label{eq:RL}
\max_{\pi_{\theta}}  \mathbb{E}_{x\sim \mathcal{D}, y\sim \pi_{\theta}(y \mid x)}\bigl[r_{\phi}(x, y)\bigr] - \beta\mathbb{D}_{\textrm{KL}}\bigl[\pi_{\theta}(y\mid x)\mid \mid \pi_\text{ref}(y\mid x)\bigr],
\end{equation}

where $\beta$ regulates the degree to which the policy $\pi_\theta$ is allowed to diverge from the reference policy $\pi_\text{ref}$, typically set as the SFT model $\pi^\text{SFT}$. In practical implementations, $\pi_\theta$ is also initialized from $\pi^\text{SFT}$. This constraint ensures that the policy remains within the region where the reward model is reliable, thus encouraging output diversity and mitigating against collapse to a small set of high-reward responses. Since language model outputs are discrete, the optimization of this objective is non-differentiable and generally tackled using reinforcement learning techniques. A standard strategy \citep{ziegler2020finetuning, stiennon2022learning, bai2022training, ouyang2022training} constructs a modified reward ${r(x, y) = r_{\phi}(x, y) -\beta (\log \pi_{\theta}(y\mid x) - \log \pi_\text{ref}(y\mid x))}$, with the maximization performed using Proximal Policy Optimization (PPO) \cite{schulman2017proximal}.

\section{Direct Preference Optimization}\label{sec:DPO}

Addressing the difficulties of using RL-based optimization for large models such as LMs, we seek a streamlined method for policy optimization that leverages human preferences directly. Contrary to previous RLHF pipelines—which first train a reward model before optimizing it with RL—our method adopts a specific reward model formulation that allows for a closed-form derivation of the corresponding optimal policy, entirely bypassing the reinforcement learning phase. As described in the following, our main innovation is to employ an explicit mapping between reward functions and their induced optimal policies, which enables the conversion of a reward-centric loss into a policy-centric loss. This transformation by change of variables allows us to dispense with explicitly fitting a standalone reward model, yet continues to optimize in accordance with established models of human preferences, such as the Bradley-Terry framework. Consequently, the

\textbf{Reformulating the DPO objective.} Consider the reinforcement learning (RL) objective in Eq.~\ref{eq:RL}, defined with a generic reward function $r$, as is customary in existing literature. Drawing upon foundational works~\citep{peters2007reinforcement, peng2019advantage, korbak2022reinforcement, go2023aligning}, it follows that the solution optimizing the KL-regularized reward maximization objective in Eq.~\ref{eq:RL} has the structure:

\begin{equation}\label{eq:op_policy}
    \pi_r(y\mid x) = \frac{1}{Z(x)}\pi_\text{ref}(y\mid x)\exp\left(\frac{1}{\beta}r(x, y)\right),
\end{equation}

where $Z(x) =\sum_{y}\pi_\text{ref}(y\mid x)\exp\left(\frac{1}{\beta}r(x, y)\right)$ serves as the normalization constant. Further mathematical details are provided in Appendix~\ref{app:derivation1}. While using an MLE-based estimate $r_\phi$ for the true reward function $r^*$, estimating the normalization term $Z(x)$ remains computationally intensive~\citep{korbak2022reinforcement, go2023aligning}, rendering this policy form less tractable for real-world implementation. Nevertheless, Eq.~\ref{eq:op_policy} can be manipulated to rewrite the reward as a function of the induced optimal policy $\pi_r$, the reference policy $\pi_\text{ref}$, and the yet-unknown normalization $Z(\cdot)$. By taking logarithms of both sides of Eq.~\ref{eq:op_policy} and rearranging terms, one obtains:

\begin{equation}\label{eq:main_eq}
    r(x,y) =\beta \log \frac{\pi_r(y\mid x)}{\pi_\text{ref}(y\mid x)} + \beta \log Z(x).
\end{equation}

This transformation applies likewise to the true reward $r^*$ and the associated optimal policy $\pi^*$. The key observation is that the Bradley-Terry model, when used for human preferences, depends solely on the difference in rewards between two outputs, i.e., ${p^*(y_1 \succ y_2 \mid x) = \sigma(r^*(x, y_1) - r^*(x, y_2))}$. Plugging the parametrization from Eq.~\ref{eq:main_eq} into the Bradley-Terry formulation in Eq.~\ref{eq:bradley-terry} for $r^*(x, y)$ leads to a cancellation of the $Z(x)$ terms. As a result, the likelihood of a human preferring $y_1$ over $y_2$ depends exclusively on the ratio of the optimal policy $\pi^*$ and the reference policy $\pi_\text{ref}$, taking the form:

\begin{equation}\label{eq:objective}
    p^*(y_1\succ y_2 \mid x)=\frac{1}{1 + \exp\left(\beta \log \frac{\pi^*(y_2\mid x)}{\pi_\text{ref}(y_2\mid x)} - \beta \log \frac{\pi^*(y_1\mid x)}{\pi_\text{ref}(y_1\mid x)}\right)}
\end{equation}

This result is elaborated in Appendix~\ref{app:derivation2}. Although Eq.~\ref{eq:objective} is expressed for the Bradley-Terry model, analogous derivations are possible for other ranking models such as Plackett-Luce~\citep{plackett1975analysis, luce2012individual}, as detailed in Appendix~\ref{app:plackett_luce_models}.

With human preference probabilities now written directly in terms of the optimal policy, it is possible to set up a maximum likelihood learning objective for a parameterized policy $\pi_\theta$. This is directly analogous to the reward modeling formulation in Eq.~\ref{eq:reward_model}, yielding the following optimization criterion:

\begin{equation}\label{eq:optimum_model}
    \mathcal{L}_\text{DPO}(\pi_{\theta}; \pi_\text{ref}) = -\mathbb{E}_{(x, y_w, y_l)\sim \mathcal{D}}\left[\log \sigma \left(\beta \log \frac{\pi_{\theta}(y_w\mid x)}{\pi_\text{ref}(y

In this approach, we employ an alternative parameterization to model an implicit reward, for which the optimal policy coincides with $\pi_\theta$. Furthermore, as the procedure essentially fits a reparametrized Bradley-Terry model, it inherits desirable theoretical properties—such as consistency—assuming certain regularity conditions on the preference data distribution \cite{bong2022generalized}. We elaborate on the theoretical aspects of DPO and how they compare to existing literature in Section~\ref{sec:theory}.

\textbf{What does the DPO update accomplish?} To gain mechanistic insights into DPO, it is helpful to study the gradient of the objective $\mathcal{L}_\text{DPO}$. Specifically, the gradient with respect to model parameters $\theta$ is expressed as:

\begin{multline*}\label{eq:gradient}
    \nabla_\theta \mathcal{L}_\text{DPO}(\pi_\theta;\pi_\text{ref}) = \\ -\beta\mathbb{E}_{(x, y_w, y_l) \sim \mathcal{D}} \bigg[\underbrace{\sigma(\hat{r}_\theta(x, y_l) - \hat{r}_\theta (x, y_w))}_\text{greater emphasis when the reward estimate misorders completions}\bigg[\underbrace{\nabla_\theta\log \pi(y_w \mid x)}_\text{pushes up probability of $y_w$} - \underbrace{\nabla_\theta\log\pi(y_l \mid x)}_\text{pushes down probability of $y_l$}\bigg]\bigg],
\end{multline*}

where $\hat{r}_\theta(x, y) = \beta \log \frac{\pi_\theta(y \mid x)}{\pi_\text{ref}(y \mid x)}$ specifies the reward implicitly characterized by the learned policy $\pi_\theta$ in relation to the reference policy $\pi_\text{ref}$ (further discussed in Section~\ref{sec:theory}). Conceptually, the loss gradient for $\mathcal{L}_\text{DPO}$ acts to elevate the likelihood of favorable completions $y_w$ while reducing the likelihood of less preferred ones $y_l$. Crucially, the weighting of each sample corresponds to the extent that the implicit reward model $\hat{r}_\theta$ assigns a higher value to a less preferred response, modulated by $\beta$. This reflects the degree to which the reward model incorrectly orders candidate outputs, in accordance with the strength of the KL regularization. Experimental findings underscore the significance of this weighting; omitting the weighting factor, as seen in a simpler variant of the method, often leads to severe degradation of model outputs (see Appendix Table~\ref{tab:unlikelihood_generations}).

\textbf{DPO Overview.} The standard workflow for DPO comprises the following steps: 1) For each prompt $x$, generate completion samples $y_1, y_2 \sim \pi_\text{ref}(\cdot \mid x)$ and annotate these samples using human preference signals to form an offline preference dataset $\mathcal{D} = \{x^{(i)}, y_w^{(i)}, y_l^{(i)}\}_{i=1}^N$, and 2) train the language model $\pi_\theta$ by minimizing $\mathcal{L}_\text{DPO}$, given the reference policy $\pi_\text{ref}$, the dataset $\mathcal{D}$, and the chosen parameter $\beta$. 

In real-world scenarios, it is often preferable to utilize pre-existing publicly released preference datasets rather than collecting new samples or annotations. Since these datasets have typically been generated using a policy $\pi^\text{SFT}$, we set $\pi_\text{ref} = \pi^\text{SFT}$ when such a model is available. If $\pi^\text{SFT}$ cannot be accessed, $\pi_\text{ref}$ is instead initialized by maximizing the likelihood over the set of preferred completions ${(x, y_w)}$, that is, we let ${\pi_\text{ref} = \argmax_{\pi}\mathbb{E}_{x, y_w \sim \mathcal{D}}\left[\log \pi(y_w \mid x)\right]}$. Adopting this approach helps reduce the mismatch between the inaccessible true reference distribution and the $\pi_\text{ref}$ applied during DPO. Additional implementation details and hyperparameter configurations are provided in Appendix~\ref{app:implementation}.

\section{Theoretical Analysis of DPO} In this section, we elaborate on the theoretical foundations of the DPO approach, offer deeper interpretations, and connect the benefits of DPO to limitations encountered in actor-critic-based RLHF algorithms (notably PPO~\cite{schulman2017proximal}).

\label{sec:theory}

\subsection{Your Language Model Is Secretly a Reward Model} DPO avoids both the explicit fitting of a reward model and the necessity of reinforcement learning to derive the optimal policy, relying solely on a maximum likelihood optimization. The objective in Eq. \ref{eq:main_eq} corresponds to a Bradley-Terry framework, where the reward function is parametrized as $r^*(x, y) = \beta \log\frac{\pi^*_\theta(y \mid x)}{\pi_\text{ref}(y \mid x)}$. The process of training the parametric model $\pi_{\theta}$ thus becomes equivalent to optimizing a reward model as in Eq. \ref{eq:reward_model} under an appropriate change of variables. Here, we develop the underlying theory behind this alternative parameterization, demonstrate that it does not impose restrictions on the possible reward models, and establish that it permits exact recovery of the optimal policy. To proceed, we introduce an equivalence relation for reward functions.

\begin{definition}
Two reward functions $r(x, y)$ and $r'(x, y)$ are defined as equivalent if and only if there exists a function $f$ such that $r(x, y)-r'(x, y) = f(x)$.
\end{definition}

This equivalence can be easily verified to satisfy the requirements of an equivalence relation, resulting in a partitioning of all reward functions into distinct equivalence classes. This gives rise to the following two lemmas:

\begin{lemma}\label{lemma:same_prefrence}
Within the Plackett-Luce preference model, including the specific case of Bradley-Terry, any pair of reward functions that are in the same equivalence class will yield identical distributions over preferences.
\end{lemma}

\begin{lemma}\label{lemma:same_policy}
    If two reward functions are members of the same equivalence class, then they produce the same optimal policy for the associated constrained reinforcement learning problem.
\end{lemma}

The proofs for these statements are direct, and are provided in Appendix \ref{app:lemma1}. The first lemma highlights the well-known issue of under-specification found in Plackett-Luce models \cite{plackett1975analysis}, whereby it is impossible to distinguish between all functions in a class without further constraints. To ensure identifiability and guarantee meaningful maximum likelihood estimates as in Eq. \ref{eq:reward_model}, it is common to impose additional assumptions or normalizations \cite{bong2022generalized}. The second lemma establishes that every reward function within an equivalence class leads to the same optimal policy; thus, our objective becomes recovering a representative reward function from the optimal class rather than any specific one. In Appendix~\ref{app:thm1}, we establish the following theorem:

\begin{theorem}\label{thm:main}
    Assuming mild technical conditions, every equivalence class of reward functions that agrees with the Plackett-Luce (and specifically Bradley-Terry) models can be parameterized as ${r(x, y) = \beta \log \frac{\pi(y\mid x)}{\pi_\text{ref}(y\mid x)}}$ for some policy $\pi(y\mid x)$, relative to a fixed reference model $\pi_\text{ref}(y \mid x)$.
\end{theorem}

\begin{sproof}
    Let $r(x, y)$ be any given reward function, associated with a unique optimal policy $\pi_r(y \mid x)$ as described in Eq. \ref{eq:op_policy}. We seek to demonstrate that a member of its equivalence class can be written in the claimed reparameterized form. Define the projection operator $f$ as  
\begin{equation}
    f(r; \pi_\text{ref}, \beta)(x, y) = r(x, y) - \beta\log\sum_{y}\pi_\text{ref}(y\mid x)\exp\left(\frac{1}{\beta}r(x, y)\right)
\end{equation}
Here, $f$ effectively centers $r(x, y)$ by subtracting the log partition function associated with $\pi_r$; since this normalization term depends only on $x$, the modified reward remains in the equivalence class of the original $r(x, y)$. Substituting the right-hand side of Eq.~\ref{eq:main_eq} for $r$, which is valid for any reward, gives $f(r; \pi_\text{ref}, \beta)(x, y) = \beta \log \frac{\pi_r(y\mid x)}{\pi_\text{ref}(y\mid x)}$. Thus, the projection $f$ provides a representative in the equivalence class with the specified reparameterization, ensuring full generality of the model.
\end{sproof}

Theorem~\ref{thm:main} can be reinterpreted as providing a precise characterization of the specific reward function that the DPO reparameterization picks from each equivalence class. Specifically, it identifies the reward function that satisfies:

\begin{equation}\label{eq:lag_p}
     \sum_{y}\underbrace{\pi_\text{ref}(y\mid x)\exp\left(\frac{1}{\beta}r(x, y)\right)}_{=\pi(y\mid x)\text{, using Thm.~\ref{thm:main} reparam.}} = 1,
\end{equation}

In other words, $\pi(y\mid x)$ is ensured to be a legitimate probability distribution—it assigns nonnegative probabilities that sum to unity. Examining Eq.~\ref{eq:op_policy}, it is evident that Eq.~\ref{eq:lag_p} corresponds to the partition function of the optimal policy determined by the reward $r(x, y)$. The DPO method exploits the idea that, by imposing suitable constraints on the broadly-defined Plackett-Luce (and, in particular, Bradley-Terry) preference models, the expressivity of the reward model is maintained, while the computation of the optimal policy in Eq.~\ref{eq:op_policy} becomes analytically manageable for any input $x$.

\subsection{Instability of Actor-Critic Algorithms} Our framework also enables the identification of instability issues present in conventional actor-critic approaches for RLHF, such as PPO. Focusing on the RL fine-tuning procedure discussed in Section~\ref{section:prelims}, we can relate our approach to the control-as-inference paradigm~\cite{levine2018reinforcement} as applied to the constrained RL setup described in \ref{eq:RL}. Assume a parametric model $\pi_{\theta}(y\mid x)$. The goal is to minimize the divergence $\mathbb{D}_{\text{KL}}[\pi_{\theta}(y|x) \mid \mid \pi^*(y\mid x)]$, where $\pi^*$ refers to the optimal policy (see Eq.~\ref{eq:optimum_model}) induced by the reward function $r_{\phi}(y, x)$. After suitable algebraic manipulations, this results in the following optimization criterion:

\begin{equation}\label{eq:AC}
    \max_{\pi_{\theta}}\mathbb{E}_{\pi_{\theta}(y\mid x)}\bigg[\underbrace{r_{\phi}(x, y) -\beta\log\sum_{y}\pi_\text{ref}(y\mid x)\exp\left(\frac{1}{\beta}r_{\phi}(x, y)\right)}_{f(r_{\phi}, \pi_\text{ref}, \beta)} - \underbrace{\beta\log\frac{\pi_{\theta}(y\mid x)}{\pi_\text{ref}(y\mid x)}}_{\text{KL}}\bigg]
\end{equation}

This objective mirrors those optimized in earlier studies 
\citep{ziegler2020finetuning, stiennon2022learning, bai2022training, ouyang2022training}, where the DPO-equivalent reward is used for the reward class $r_{\phi}$. In this framework, the normalization factor within $f(r_{\phi}, \pi_\text{ref}, \beta)$ can be viewed as the soft value function corresponding to the reference policy $\pi_\text{ref}$. Although this normalization does not influence the optimal policy, omitting it can lead to a policy gradient with substantial variance, thereby destabilizing the learning process. Employing a learned value function to approximate the normalization is possible, but it poses its own optimization challenges. Previous approaches often normalize rewards by leveraging a human completion as a baseline, which amounts to a single-sample Monte-Carlo estimate of the normalization factor. In contrast, the DPO reparameterization produces a reward function inherently free from any reliance on baseline estimates.

\section{Experiments}
Here, we present an empirical assessment of DPO for direct policy optimization from preference data. Initially, within a controlled text-generation scenario, we investigate the efficiency with which DPO balances the dual objectives of reward maximization and minimizing the KL-divergence from the reference policy, benchmarking its performance against established preference optimization algorithms such as PPO. Subsequently, we extend our analysis to encompass larger-scale models and more challenging RLHF problems, including tasks in summarization and conversational dialogue. Our findings reveal that, with minimal hyperparameter tuning, DPO often achieves parity with or outperforms competitive baselines, such as RLHF with PPO and strategies that select the top result from $N$ sampled trajectories according to a learned reward model. Prior to reporting these outcomes, we provide an overview of our experimental protocol, with further specifics available in Appendix~\ref{app:exp_details}.

\textbf{Tasks.} Our study investigates three distinct open-ended text generation tasks. For each task, the various algorithms optimize a policy using a dataset of preferences denoted by $\mathcal{D}=\bigl\{x^{(i)}, y_w^{(i)}, y_l^{(i)}\bigr\}_{i=1}^N$. In the case of \textbf{controlled sentiment generation}, $x$ represents the beginning of a movie review sourced from the IMDb dataset \cite{maas-EtAl:2011:ACL-HLT2011}, and the objective of the policy is to produce $y$ with a positive sentiment orientation. To facilitate controlled evaluation in this setting, we construct preference pairs automatically for each generated output by utilizing a pre-trained sentiment classifier, ensuring that $p(\text{positive}\mid x,y_w)>p(\text{positive}\mid x,y_l)$. For SFT, the GPT-2-large model is fine-tuned to convergence using the training partition of the IMDb reviews (additional details in App~\ref{app:sentiment_details}). For the \textbf{summarization} task, $x$ is drawn from Reddit forum posts, and the policy is required to generate summaries $y$ capturing the primary points of each post. In line with earlier research, we employ the Reddit TL;DR summarization dataset \citep{volske-etal-2017-tl} and utilize human preferences collected by \citeauthor{stiennon2022learning}. SFT models are further fine-tuned on human-authored summaries of Reddit posts\footnote{\url{https://huggingface.co/CarperAI/openai_summarize_tldr_sft}} using the TRLX framework for RLHF \citep{leandro_von_werra_2023_7790115}. Human preference annotations are derived from samples generated by a model trained in a manner similar, but not identical, to our SFT setup, as reported by \citeauthor{stiennon2022learning}. For the \textbf{single-turn dialogue} task, $x$ is a prompt from a human user, spanning a wide array of queries, including those about astrophysics or requests for personal advice. The policy must deliver a relevant and constructive reply $y$; this task uses the Anthropic Helpful and Harmless dialogue dataset \citep{bai2022training}, comprising 170,000 dialogue episodes between human users and a conversational agent. Each dialogue ends with two model-generated responses, accompanied by human judgments indicating which is preferred. In this domain, a pre-trained SFT model is not available, so we instead construct the SFT model by fine-tuning a publicly available language model exclusively on human-preferred completions.

\textbf{Evaluation.} We adopt two complementary evaluation strategies in our experimental framework. For the controlled sentiment generation scenario, where the ground-truth reward function (specifically, a sentiment classifier) is available, each method is assessed by mapping the trade-off frontier between achieved reward and KL-divergence from the reference policy; this is feasible due to direct access to the true reward values. Conversely, in real-world contexts where the ground-truth reward function remains unavailable, we rely on the \textit{win rate} metric, benchmarking each method against a baseline policy. To serve as a stand-in for human evaluators, we employ GPT-4, which scores summary quality in summarization tasks and evaluates response helpfulness for single-turn dialogue. For summarization tasks, the test set’s reference summaries serve as the baseline comparator, while for dialogue tasks, we use the annotator-preferred responses from the dataset as the baseline. While prior work suggests that large language models are often superior to traditional automatic metrics as evaluators \citep{Chen2023ExploringTU}, we also conduct a human user study (see Sec.~\ref{sec:human-judgments}) to corroborate our reliance on GPT-4. Our findings indicate a strong alignment between GPT-4's ratings and human judgments, with the agreement between humans and GPT-4 equaling or surpassing typical inter-annotator concordance.

\textbf{Methods.} Beyond DPO, we compare a variety of established training regimes for language models aimed at following human preferences. For the summarization task, we utilize zero-shot prompting with \textbf{GPT-J} \citep{gpt-j}, and for the dialogue task, we apply 2-shot prompting to \textbf{Pythia-2.8B} \citep{biderman2023pythia}. Furthermore, we test the \textbf{SFT} model, and the \textbf{Preferred-FT} variant—fine-tuned via supervised learning specifically on the selected completions $y_w$, sourced either from SFT outputs (in sentiment and summarization settings) or from a general language model (for dialogue tasks). Among pseudo-supervised methods, we include \textbf{Unlikelihood}~\citep{welleck2019neural}, which trains the policy to maximize the likelihood of $y_w$ and explicitly minimize the probability assigned to $y_l$; this is regulated by an optional `unlikelihood' weight $\alpha \in [0, 1]$. We also incorporate \textbf{PPO} \citep{schulman2017proximal}, using a reward model trained on preference data, and its oracle variant \textbf{PPO-GT}, which leverages the actual reward function in the controlled sentiment environment. In sentiment experiments, PPO-GT is instantiated using two approaches: an out-of-the-box implementation \cite{leandro_von_werra_2023_7790115}, and an enhanced version featuring reward normalization and hyperparameter refinement for improved outcomes; similar adjustments are employed in PPO with learned rewards. Lastly, we evaluate a \textbf{Best of $N$} baseline, which entails generating $N$ completions from the SFT model (or Preferred-FT in dialogue) and selecting the top-scoring output per a reward model derived from preference data. While this approach is typically strong, its practicality is hindered at even moderate values of $N$, due to the computational cost of sampling $N$ responses for each test query.

\subsection{How well can DPO optimize the RLHF objective?}

The reward maximization objective employed in standard RLHF approaches introduces a constraint on the KL divergence, balancing the desire to increase reward with the need to limit the policy’s divergence from the reference policy. As such, algorithm comparisons should jointly consider both the obtained reward and the corresponding KL divergence; a method that secures marginally greater reward at the cost of a substantially increased KL is not inherently preferable. In Figure~\ref{fig:frontier-tldr-main}, we display the tradeoff between reward and KL for several algorithms within the sentiment task. Each algorithm is trained multiple times, each instance using a distinct value for the policy conservativeness parameter (target KL $\in\{3,6,9,12\}$ for PPO, $\beta \in \{0.05,0.1,1,5\}$, $\alpha\in\{0.05,0.1,0.5,1\}$ for unlikelihood, and various random seeds for preferred-FT), yielding a total of 22 training runs. At intervals of 100 training steps up to convergence, policies are evaluated on a standardized set of test prompts. During this evaluation, we record both the mean reward according to the actual reward function and the mean sequence-level KL\footnote{Meaning, the aggregate of KL-divergences computed at each timestep.} with respect to the reference policy $\text{KL}\left(\pi\mid \mid \pi_\text{ref}\right)$. The findings indicate that DPO constructs a significantly more favorable frontier, attaining superior reward levels without incurring high KL divergence. This finding is noteworthy for a number of reasons. Primarily, even though both DPO and PPO are designed to maximize the same objective, DPO does so with markedly greater efficiency—its reward and KL balance consistently outperforms PPO. Additionally, DPO produces a superior frontier in comparison to PPO, \emph{even when PPO is granted access to ground truth rewards} (PPO-GT).

\subsection{Can DPO scale to real preference datasets?}
\label{sec:dpo-real-datasets}
We proceed to assess the fine-tuning capabilities of DPO on tasks involving summarization and single-turn dialogue. In summarization, established metrics like ROUGE have been shown to have a weak association with human judgments~\citep{stiennon2022learning}. Prior research suggests that language models fine-tuned with PPO on human preference data tend to yield more favorable summaries. Our evaluation entails generating model outputs using the test portion of the TL;DR summarization dataset and measuring each method's mean win rate relative to the reference summaries in this test split. For all approaches, we sample outputs across a temperature range from 0.0 to 1.0 and present the corresponding win rates in Figure~\ref{fig:frontier-tldr-main} (right). The models for DPO, PPO, and Preferred-FT are each initialized from the same GPT-J SFT model\footnote{\url{https://huggingface.co/CarperAI/openai_summarize_tldr_sft}}. Experimental results reveal that DPO attains a win rate close to 61\% at a sampling temperature of 0.0, outperforming PPO, which achieves around 57\% at its optimal sampling temperature (0.0). Additionally, DPO secures a higher peak win rate than the best-of-$N$ baseline. It should be emphasized that DPO's $\beta$ hyperparameter was not extensively tuned, implying that actual performance might be understated. Furthermore, DPO demonstrates notably greater resilience to changes in sampling temperature compared to PPO, whose performance can regress to that of the original GPT-J model at higher temperatures. The Preferred-FT method yields only marginal gains relative to the SFT baseline. We further examine DPO and PPO through human preference tests in Section~\ref{sec:human-judgments}, observing that DPO completions sampled at temperature 0.25 were favored in 58\% of cases over PPO completions at the same temperature.

For the single-turn dialogue evaluation, we assess various approaches using the portion of the Anthropic HH dataset’s test split \citep{bai2022training} that consists of one human-assistant exchange. When using GPT-4 for evaluation, we treat the preferred completions from the test set as ground truth references and calculate the win rate for each approach relative to these. In the absence of a standard SFT baseline for this scenario, our methodology begins with a pre-trained Pythia-2.8B model. We then utilize Preferred-FT to fine-tune a reference model on the preferred completions, ensuring outputs remain within the data distribution, followed by DPO training. Additionally, we benchmark against two baselines: the best-performing completion among 128 Preferred-FT samples (noting that increasing beyond 128 yields diminishing returns—see Appendix Figure~\ref{fig:best-of-n}) and a 2-shot prompted version of the Pythia-2.8B base model. Our findings indicate that DPO matches or surpasses the other methods at their respective optimal temperature settings. We further assess a PPO-trained RLHF model on the Anthropic HH dataset\footnote{\url{https://huggingface.co/reciprocate/ppo_hh_pythia-6B}}, developed by a recognized source\footnote{\url{https://github.com/CarperAI/trlx/tree/main/examples/hh}}. However, we could not identify any prompting strategy or sampling temperature that allowed this PPO model to outperform the base Pythia-2.8B. Considering the parallel reward functions and supporting evidence from TL;DR experiments, we treat the Best of 128 benchmark as an approximate stand-in for PPO-level outcomes. Taken together, DPO stands out as the only approach that is computationally efficient while exceeding the performance of the Anthropic HH dataset’s preferred completions, and it delivers results on par with or better than the much more resource-intensive Best of 128 standard. Lastly, Figure~\ref{fig:dialogue-main} demonstrates that DPO attains its peak performance rapidly.

\subsection{Generalization to a new input distribution}

To further analyze the robustness of PPO and DPO in the face of distributional shifts, we apply the policies trained on Reddit TL;DR summarization to a distinctly different domain: news articles from the CNN/DailyMail test set \citep{nallapati-etal-2016-abstractive}. These evaluations employ the optimal sampling temperatures determined for TL;DR, specifically 0 and 0.25. Table~\ref{tab:ood} displays the outcomes. We calculated the GPT-4 win rate against the ground-truth summaries of this dataset using the same GPT-4 (C) prompt as in the Reddit TL;DR analysis, with the phrase ``forum post'' replaced by ``news article.'' On this out-of-distribution benchmark, DPO continues to substantially outperform the PPO policy. These findings provide preliminary support that DPO policies generalize comparably to those trained with PPO, despite DPO’s lack of exposure to the extra unlabeled Reddit TL;DR prompts utilized by PPO.

\subsection{Validating GPT-4 judgments with human judgments}
\label{sec:human-judgments}
A human subject experiment was conducted to assess the trustworthiness of GPT-4's evaluations, leveraging results from the TL;DR summarization study and two alternative prompts for GPT-4. The \textbf{GPT-4 (S)} (simple) prompt requests a judgment as to which summary better encapsulates the critical information in a post. In contrast, the \textbf{GPT-4 (C)} (concise) prompt additionally requests an assessment of conciseness, motivated by observations that the \textbf{GPT-4 (S)} prompt often leads GPT-4 to favor longer and more redundant summaries relative to human preferences. Full prompt texts are included in Appendix~\ref{app:prompts}. We examine three model configurations: the highest scoring (DPO, temp. 0.25), the lowest scoring (PPO, temp. 1.0), and a representative mid-range performer (SFT, temp. 0.25), in order to capture a broad spectrum of sample quality; all three are evaluated against the PPO policy at its optimal (greedy) sampling temperature. Our findings indicate that for both prompt variants, GPT-4’s agreement rate with humans closely matches human inter-annotator agreement rates, implying that GPT-4 serves as a credible proxy for human evaluation (noting that multiple human ratings were collected only for the DPO and PPO-1 pairings due to limited rater availability). On balance, the \textbf{GPT-4 (C)} prompt yields win rates that more closely mirror human judgments; hence, this prompt is used for the primary results reported in Section~\ref{sec:dpo-real-datasets}. Further information on the human subject evaluation, including details on the web-based rater interface and the composition of the volunteer pool, is provided in Appendix~\ref{app:human-study}.

\section{Discussion}
Preference-based learning offers a robust and scalable approach for developing capable and well-aligned language models. In this work, we presented DPO, a straightforward framework that enables training language models from preferences without relying on reinforcement learning. Instead of reshaping preference learning into a conventional RL framework to leverage standard RL algorithms, DPO establishes a correspondence between language model policies and reward functions, facilitating the direct optimization of models towards human preferences via a simple cross-entropy loss function—eschewing reinforcement learning without sacrificing generality. DPO achieves performance on par with or exceeding leading RLHF methods, such as those employing PPO, often with minimal need for hyperparameter adjustment; as a result, DPO significantly streamlines the process of training language models based on human feedback.

\textbf{Limitations \& Future Work.} Our findings open several avenues for subsequent investigation. For instance, how well does a policy trained via DPO generalize to data distributions distinct from its training set, especially in comparison to approaches that utilize explicit reward models? Preliminary results indicate that DPO-trained models may generalize comparably to those trained with PPO, but further systematic study is warranted. Additionally, it remains an open question whether approaches such as self-labeling with the DPO policy can make similarly effective use of prompts that lack explicit annotations. Another topic for future study is understanding the nature of reward over-optimization within the DPO context and whether the modest dip in performance observed in Figure~\ref{fig:dialogue-main}-right might exemplify this phenomenon. Furthermore, although our evaluation is limited to models up to 6B parameters, future research should explore the feasibility and outcomes of applying DPO at scales comparable to cutting-edge models with significantly larger capacities. In terms of assessment methodology, our experiments suggest that GPT-4’s calculated win rates are influenced by prompt phrasing; identifying optimal protocols for obtaining reliable judgments from automated evaluators remains a promising research challenge. Lastly, DPO’s utility is not confined to language models—extending its application to other generative modeling domains offers intriguing possibilities for future exploration.